\begin{lexlean}{Main}
\useglossary{lexlean.std.nat@1.0.0}
\useglossary{peano.arith@1.0.0}
\importmodule{Addition}
\importmodule{Divisibility}
\importmodule{Doubling}
\importmodule{Order}
\title{Natural number addition: consequences}

\begin{section}{consequences}
\heading{Commutativity: consequences}
  \begin{theorem}{comm-under-hypothesis}
  \noaxioms
  For every natural number \(x\) and natural number \(y\), if \(x = y\), then \(x + y = y + x\).
  \begin{proof}
  Assume \(h\).
  Close the goal with \(\reference{Addition::add-comm}(x, y)\).
  \end{proof}
  \end{theorem}

  \begin{theorem}{comm-rewrite}
  \noaxioms
  For every natural number \(x\) and natural number \(y\), if \(x = y\), then \(x + x = y + x\).
  \begin{proof}
  Assume \(h\).
  \begin{rewrite}{goal}
  \forward{h}
  \end{rewrite}
  \end{proof}
  \end{theorem}

  \begin{theorem}{double-divides}
  \noaxioms
  For every natural number \(n\), \(2\) divides \(double(n)\).
  \begin{proof}
  Use \(n\) as the witness.
  Close the goal with \(eqsymm(twomul(n))\).
  \end{proof}
  \end{theorem}

  \begin{theorem}{even-divides}
  \noaxioms
  For every natural number \(n\), if \(n\) is even, then \(2\) divides \(n\).
  \begin{proof}
  Assume \(h\).
  \begin{cases}{h}
  \begin{case}{lexlean.core::exists-intro}
  \bind{k;hk}
  Use \(k\) as the witness.
  \begin{rewrite}{goal}
  \forward{hk}
  \end{rewrite}
  Close the goal with \(eqsymm(twomul(k))\).
  \end{case}
  \end{cases}
  \end{proof}
  \end{theorem}
\end{section}

\begin{section}{excluded-middle}
\heading{Excluded middle}
  \begin{theorem}{even-or-not}
  \exactaxioms{Classical.choice;Quot.sound;propext}
  For every natural number \(n\), \(n\) is even or not \(n\) is even.
  \begin{proof}
  Close the goal with \(em(\lexeme{peano.arith::even}(n))\).
  \end{proof}
  \end{theorem}

  \begin{theorem}{positive-or-zero}
  \allowaxioms{Classical.choice;Quot.sound;propext}
  For every natural number \(n\), \(n\) is positive or \(n = 0\).
  \begin{proof}
  \begin{cases}{n}
  \begin{case}{lexlean.std.nat::zero}
  \bind{}
  Select the right alternative.
  Close the goal by reflexivity.
  \end{case}
  \begin{case}{lexlean.std.nat::succ}
  \bind{m}
  Select the left alternative.
  Close the goal with \(\reference{Order::successor-positive}(m)\).
  \end{case}
  \end{cases}
  \end{proof}
  \end{theorem}
\end{section}
\end{lexlean}
